\conclusion

В ходе выполнения выпускной квалификационной работы была разработана система для сравнительного нагрузочного тестирования и визуального анализа производительности реляционных баз данных. Работа выполнена в полном объёме, все поставленные задачи решены.

В ходе аналитического исследования проведён обзор современных подходов к нагрузочному тестированию СУБД, рассмотрены методологии стандартизированного тестирования TPC-C, TPC-H и TPC-E, выполнен сравнительный анализ шести инструментов (HammerDB, Sysbench, pgbench, Oracle RAT, BenchmarkSQL, YCSB). Выявлено, что ни одно из существующих решений не обеспечивает полный цикл нагрузочного тестирования с визуальным анализом результатов, хранением истории и автоматизированным сравнением прогонов в едином веб-интерфейсе. Сформулированы функциональные и нефункциональные требования к разрабатываемой системе.

Спроектирована многоуровневая клиент-серверная архитектура системы с разделением на серверную часть (FastAPI) и клиентскую часть (Next.js). Спроектирован модуль генерации нагрузки на основе асинхронного выполнения SQL-запросов с параллельной обработкой виртуальных пользователей, сбором системных и внутренних метрик СУБД, механизмом параметризации запросов и фазой прогрева. Спроектирован модуль управления состоянием базы данных, обеспечивающий автоматическое создание резервных копий перед тестами с операциями записи и восстановление после завершения теста с верификацией целостности. Спроектирована подсистема визуализации и сравнительного анализа с применением статистических критериев (Welch's t-test, Mann-Whitney U). Разработана модель данных, включающая сущности TestRun, TestResult, TimeSeries, MetricSample, TestScenario и DatabaseConnectionConfig.

Реализована серверная часть системы: модуль генерации нагрузки, модуль управления состоянием БД, модуль сравнительного анализа, REST API с семью группами эндпоинтов, WebSocket-менеджер для трансляции метрик в реальном времени. Реализована клиентская часть: Zustand store для управления состоянием, API-клиент с типизированными методами, WebSocket-хук для получения метрик в реальном времени, семь страниц интерфейса с визуализацией через Recharts. Развёртывание организовано через Docker Compose для базы данных истории и отдельные контейнеры для тестовых СУБД.

Выполнено тестирование системы на тестовых стендах: MySQL с тестовой схемой Sakila и PostgreSQL с тестовой схемой Pagila. Продемонстрировано функционирование всех компонентов системы: создание подключений, конфигурирование и запуск тестов, мониторинг метрик в реальном времени, просмотр истории и сравнительный анализ. Подтверждено соответствие системы всем сформулированным функциональным и нефункциональным требованиям.

Результаты работы могут быть использованы инженерами по производительности, DevOps-инженерами, администраторами баз данных и разработчиками для обоснованного выбора СУБД, оценки эффективности оптимизаций и предотвращения регрессии производительности. Система может быть расширена путём добавления поддержки дополнительных СУБД, реализации горизонтального масштабирования генераторов нагрузки и интеграции с системами непрерывной интеграции.
